\documentclass[10pt,a4paper]{amsart}
\usepackage{amsmath,amssymb,amsthm,graphicx}
% Retained preamble for the Stacks Project chapter snapshot in this directory.
% The chapter's own preamble.tex is a copy of the shared Stacks Project preamble; the block
% below is sliced MECHANICALLY, line for line, from it. Nothing was added, removed or
% redefined; the wrapper class, the xr-hyper/hyperref wiring and the externaldocument list
% are omitted because the assembler supplies the document class.
%
% lines taken from preamble.tex: 3, 5, 9, 11, 12, 156, 157, 158, ... (51 lines total)
%
% package coverage restored from the chapter's own preamble (2026-10-13): xy, xy-2cell

\usepackage[all]{xy}
\xyoption{2cell}
\UseAllTwocells

\newenvironment{reference}{\comment}{\endcomment}
\newenvironment{slogan}{\comment}{\endcomment}
\newenvironment{history}{\comment}{\endcomment}
\theoremstyle{plain}
\newtheorem{theorem}[subsection]{Theorem}
\newtheorem{proposition}[subsection]{Proposition}
\newtheorem{lemma}[subsection]{Lemma}
\theoremstyle{definition}
\newtheorem{definition}[subsection]{Definition}
\newtheorem{example}[subsection]{Example}
\newtheorem{exercise}[subsection]{Exercise}
\newtheorem{situation}[subsection]{Situation}
\theoremstyle{remark}
\newtheorem{remark}[subsection]{Remark}
\newtheorem{remarks}[subsection]{Remarks}
\numberwithin{equation}{subsection}
\def\lim{\mathop{\mathrm{lim}}\nolimits}
\def\colim{\mathop{\mathrm{colim}}\nolimits}
\def\Spec{\mathop{\mathrm{Spec}}}
\def\Hom{\mathop{\mathrm{Hom}}\nolimits}
\def\Ext{\mathop{\mathrm{Ext}}\nolimits}
\def\SheafHom{\mathop{\mathcal{H}\!\mathit{om}}\nolimits}
\def\SheafExt{\mathop{\mathcal{E}\!\mathit{xt}}\nolimits}
\def\Sch{\mathit{Sch}}
\def\Mor{\mathop{\mathrm{Mor}}\nolimits}
\def\Ob{\mathop{\mathrm{Ob}}\nolimits}
\def\Sh{\mathop{\mathit{Sh}}\nolimits}
\def\NL{\mathop{N\!L}\nolimits}
\def\CH{\mathop{\mathrm{CH}}\nolimits}
\def\proetale{{pro\text{-}\acute{e}tale}}
\def\etale{{\acute{e}tale}}
\def\QCoh{\mathit{QCoh}}
\def\Ker{\mathop{\mathrm{Ker}}}
\def\Im{\mathop{\mathrm{Im}}}
\def\Coker{\mathop{\mathrm{Coker}}}
\def\Coim{\mathop{\mathrm{Coim}}}
\DeclareMathSymbol{\boxtimes}{\mathbin}{AMSa}{"02}
\def\QCohstack{\mathcal{QC}\!\mathit{oh}}
\def\Cohstack{\mathcal{C}\!\mathit{oh}}
\def\Spacesstack{\mathcal{S}\!\mathit{paces}}
\def\Quotfunctor{\mathrm{Quot}}
\def\Hilbfunctor{\mathrm{Hilb}}
\def\Curvesstack{\mathcal{C}\!\mathit{urves}}
\def\Polarizedstack{\mathcal{P}\!\mathit{olarized}}
\def\Complexesstack{\mathcal{C}\!\mathit{omplexes}}
\def\Pic{\mathop{\mathrm{Pic}}\nolimits}
\def\Picardstack{\mathcal{P}\!\mathit{ic}}
\def\Picardfunctor{\mathrm{Pic}}
\def\Deformationcategory{\mathcal{D}\!\mathit{ef}}
\usepackage{cleveref}
\crefname{originallemma}{original lemma}{original lemmas}
\Crefname{originallemma}{Original lemma}{Original lemmas}
\makeatletter
\newlabel{equation-pushforward}{{equation-pushforward}{source}{}{sourcepointer-1}{}}
\newlabel{equation-pushforward@cref}{{[originallemma][0][]equation-pushforward}{[1][0][]source}{}{}{}}
\makeatother
\crefname{originallemma}{original lemma}{original lemmas}
\Crefname{originallemma}{Original lemma}{Original lemmas}
\makeatletter
\newlabel{lemma-compare}{{lemma-compare}{source}{}{sourcepointer-2}{}}
\newlabel{lemma-compare@cref}{{[originallemma][0][]lemma-compare}{[1][0][]source}{}{}{}}
\makeatother
\crefname{originallemma}{original lemma}{original lemmas}
\Crefname{originallemma}{Original lemma}{Original lemmas}
\makeatletter
\newlabel{section-presheaves}{{section-presheaves}{source}{}{sourcepointer-3}{}}
\newlabel{section-presheaves@cref}{{[originallemma][0][]section-presheaves}{[1][0][]source}{}{}{}}
\makeatother
\crefname{originallemma}{original lemma}{original lemmas}
\Crefname{originallemma}{Original lemma}{Original lemmas}
\makeatletter
\newlabel{spaces-properties-lemma-functoriality-etale-site}{{03G2}{source}{}{sourcepointer-4}{}}
\newlabel{spaces-properties-lemma-functoriality-etale-site@cref}{{[originallemma][0][]03G2}{[1][0][]source}{}{}{}}
\makeatother
\crefname{originallemma}{original lemma}{original lemmas}
\Crefname{originallemma}{Original lemma}{Original lemmas}
\makeatletter
\newlabel{spaces-properties-remark-explain-equivalence}{{03H7}{source}{}{sourcepointer-5}{}}
\newlabel{spaces-properties-remark-explain-equivalence@cref}{{[originallemma][0][]03H7}{[1][0][]source}{}{}{}}
\makeatother
\crefname{originallemma}{original lemma}{original lemmas}
\Crefname{originallemma}{Original lemma}{Original lemmas}
\makeatletter
\newlabel{spaces-properties-section-etale-site}{{03EB}{source}{}{sourcepointer-6}{}}
\newlabel{spaces-properties-section-etale-site@cref}{{[originallemma][0][]03EB}{[1][0][]source}{}{}{}}
\makeatother
\crefname{originallemma}{original lemma}{original lemmas}
\Crefname{originallemma}{Original lemma}{Original lemmas}
\makeatletter
\newlabel{topologies-lemma-morphism-big-small-etale}{{021I}{source}{}{sourcepointer-7}{}}
\newlabel{topologies-lemma-morphism-big-small-etale@cref}{{[originallemma][0][]021I}{[1][0][]source}{}{}{}}
\makeatother
\title{Stacks, Sheaves on Algebraic Stacks: lemma-compare-morphism}
\author{The Stacks Project Authors}
\date{stacks-project:stacks-sheaves.tex:local-snapshot-20261012; GNU FDL 1.2 or later}
\begin{document}
\maketitle
\noindent\textit{Editorial note.} Statement span is lines 1552-1570 of sources/stacks-algebraic-stacks-sheaves/pinned-original-source.tex; proof span is lines 1572-1668, the complete author proof (95 body lines).\par\medskip
\noindent\textit{Editorial note.} Navigation stubs: 7 cited-but-not-included labels are supplied as non-load-bearing pointers with their pinned tags where the snapshot tag index has them.\par\medskip
\noindent\textit{Editorial note.} extra\_packages is empty on purpose: the retained preamble already loads xy with options and a second declaration clashes in the assembler wrapper.\par\medskip
\noindent\textit{Editorial note.} Text is retained verbatim, including the author's own typos and placeholders.\par\medskip
\noindent\textit{Editorial note.} Queue index 4 of queue\_algebraic-stacks-sheaves.json (pb=3729, stmt 1552-1570, proof 1572-1668).\par\medskip
\noindent\textit{Allowed standard prerequisites (editorial).}\begin{itemize}
\item Explicit prerequisite boundary. The cited inputs are prerequisites; the argument in the retained span is not replaced by any cited result.
\end{itemize}\medskip
\ifdefined\hypertarget\hypertarget{sourcepointer-1}{}\fi
\noindent\textit{Source reference.} Original lemma equation-pushforward: sources/stacks-algebraic-stacks-sheaves/pinned-original-source.tex, source lines 747--751. Cited by the retained statement or proof but not part of the collected spans; supplied only as a navigation pointer. The argument does not depend on its content.\par\medskip
\ifdefined\hypertarget\hypertarget{sourcepointer-2}{}\fi
\noindent\textit{Source reference.} Original lemma lemma-compare: sources/stacks-algebraic-stacks-sheaves/pinned-original-source.tex, source lines 1430--1452. Cited by the retained statement or proof but not part of the collected spans; supplied only as a navigation pointer. The argument does not depend on its content.\par\medskip
\ifdefined\hypertarget\hypertarget{sourcepointer-3}{}\fi
\noindent\textit{Source reference.} Original lemma section-presheaves: sources/stacks-algebraic-stacks-sheaves/pinned-original-source.tex, source lines 302--319. Cited by the retained statement or proof but not part of the collected spans; supplied only as a navigation pointer. The argument does not depend on its content.\par\medskip
\ifdefined\hypertarget\hypertarget{sourcepointer-4}{}\fi
\noindent\textit{Source reference.} Original lemma 03G2: sources/stacks-algebraic-stacks-sheaves/pinned-original-source.tex, source lines 1668--1668. Cited by the retained statement or proof but not part of the collected spans; supplied only as a navigation pointer. The argument does not depend on its content.\par\medskip
\ifdefined\hypertarget\hypertarget{sourcepointer-5}{}\fi
\noindent\textit{Source reference.} Original lemma 03H7: sources/stacks-algebraic-stacks-sheaves/pinned-original-source.tex, source lines 1668--1668. Cited by the retained statement or proof but not part of the collected spans; supplied only as a navigation pointer. The argument does not depend on its content.\par\medskip
\ifdefined\hypertarget\hypertarget{sourcepointer-6}{}\fi
\noindent\textit{Source reference.} Original lemma 03EB: sources/stacks-algebraic-stacks-sheaves/pinned-original-source.tex, source lines 1668--1668. Cited by the retained statement or proof but not part of the collected spans; supplied only as a navigation pointer. The argument does not depend on its content.\par\medskip
\ifdefined\hypertarget\hypertarget{sourcepointer-7}{}\fi
\noindent\textit{Source reference.} Original lemma 021I: sources/stacks-algebraic-stacks-sheaves/pinned-original-source.tex, source lines 1668--1668. Cited by the retained statement or proof but not part of the collected spans; supplied only as a navigation pointer. The argument does not depend on its content.\par\medskip
\begin{lemma}
\label{lemma-compare-morphism}
Let $S$ be a scheme. Let $f : \mathcal{X} \to \mathcal{Y}$ be a morphism
of categories fibred in groupoids over $(\Sch/S)_{fppf}$. Assume
$\mathcal{X}$, $\mathcal{Y}$ are representable by algebraic spaces $F$, $G$.
Denote $f : F \to G$ the induced morphism of algebraic spaces, and
$f_{small} : F_\etale \to G_\etale$
the corresponding morphism of ringed topoi. Then
$$
\xymatrix{
(\Sh(\mathcal{X}_\etale), \mathcal{O}_\mathcal{X})
\ar[d]_{\pi_F} \ar[rr]_f & &
(\Sh(\mathcal{Y}_\etale), \mathcal{O}_\mathcal{Y}) \ar[d]^{\pi_G} \\
(\Sh(F_\etale), \mathcal{O}_F) \ar[rr]^{f_{small}} & &
(\Sh(G_\etale), \mathcal{O}_G)
}
$$
is a commutative diagram of ringed topoi.
\end{lemma}

\begin{proof}
This is similar to
Topologies, Lemma \ref{topologies-lemma-morphism-big-small-etale} (3)
but there is a small snag due to the fact that $F \to G$ may not be
representable by schemes. In particular we don't get a commutative diagram
of ringed sites, but only a commutative diagram of ringed topoi.

\medskip\noindent
Before we start the proof proper, we choose equivalences
$j : \mathcal{S}_F \to \mathcal{X}$ and
$j' : \mathcal{S}_G \to \mathcal{Y}$ which induce functors
$u : F_\etale \to \mathcal{X}$ and
$u' : G_\etale \to \mathcal{Y}$ as in the proof of
Lemma \ref{lemma-compare}. Because of the 2-functoriality of
sheaves on categories fibred in groupoids over $\Sch_{fppf}$
(see discussion in Section \ref{section-presheaves})
we may assume that $\mathcal{X} = \mathcal{S}_F$ and
$\mathcal{Y} = \mathcal{S}_G$ and that $f : \mathcal{S}_F \to \mathcal{S}_G$
is the functor associated to the morphism $f : F \to G$. Correspondingly
we will omit $u$ and $u'$ from the notation, i.e., given an object
$U \to F$ of $F_\etale$ we denote $U/F$
the corresponding object of $\mathcal{X}$. Similarly for $G$.

\medskip\noindent
Let $\mathcal{G}$ be a sheaf on $\mathcal{X}_\etale$.
To prove (2) we compute $\pi_{G, *}f_*\mathcal{G}$ and
$f_{small, *}\pi_{F, *}\mathcal{G}$. To do this let $V \to G$ be an object
of $G_\etale$. Then
$$
\pi_{G, *}f_*\mathcal{G}(V) = f_*\mathcal{G}(V/G) =
\Gamma\Big(
(\Sch/V)_{fppf} \times_{\mathcal{Y}} \mathcal{X},
\ \text{pr}^{-1}\mathcal{G}\Big)
$$
see (\ref{equation-pushforward}). The fibre product in the formula is
$$
(\Sch/V)_{fppf} \times_{\mathcal{Y}} \mathcal{X} =
(\Sch/V)_{fppf} \times_{\mathcal{S}_G} \mathcal{S}_F =
\mathcal{S}_{V \times_G F}
$$
i.e., it is the split category fibred in groupoids associated to the
algebraic space $V \times_G F$. And $\text{pr}^{-1}\mathcal{G}$ is a
sheaf on $\mathcal{S}_{V \times_G F}$ for the \'etale topology.

\medskip\noindent
In particular, if $V \times_G F$ is representable, i.e., if it is a scheme,
then $\pi_{G, *}f_*\mathcal{G}(V) = \mathcal{G}(V \times_G F/F)$ and
also
$$
f_{small, *}\pi_{F, *}\mathcal{G}(V) =
\pi_{F, *}\mathcal{G}(V \times_G F) =
\mathcal{G}(V \times_G F/F)
$$
which proves the desired equality in this special case.

\medskip\noindent
In general, choose a scheme $U$ and a surjective \'etale morphism
$U \to V \times_G F$. Set $R = U \times_{V \times_G F} U$. Then
$U/V \times_G F$ and $R/V \times_G F$ are objects of the fibre
product category above. Since $\text{pr}^{-1}\mathcal{G}$ is a
sheaf for the \'etale topology on $\mathcal{S}_{V \times_G F}$
the diagram
$$
\xymatrix{
\Gamma\Big(
(\Sch/V)_{fppf} \times_{\mathcal{Y}} \mathcal{X},
\ \text{pr}^{-1}\mathcal{G}\Big)
\ar[r] &
\text{pr}^{-1}\mathcal{G}(U/V \times_G F) \ar@<1ex>[r] \ar@<-1ex>[r] &
\text{pr}^{-1}\mathcal{G}(R/V \times_G F)
}
$$
is an equalizer diagram. Note that
$\text{pr}^{-1}\mathcal{G}(U/V \times_G F) = \mathcal{G}(U/F)$
and $\text{pr}^{-1}\mathcal{G}(R/V \times_G F) = \mathcal{G}(R/F)$
by the definition of pullbacks. Moreover, by the material in
Properties of Spaces, Section \ref{spaces-properties-section-etale-site}
(especially,
Properties of Spaces,
Remark \ref{spaces-properties-remark-explain-equivalence} and
Lemma \ref{spaces-properties-lemma-functoriality-etale-site})
we see that there is an equalizer diagram
$$
\xymatrix{
f_{small, *}\pi_{F, *}\mathcal{G}(V)
\ar[r] &
\pi_{F, *}\mathcal{G}(U/F) \ar@<1ex>[r] \ar@<-1ex>[r] &
\pi_{F, *}\mathcal{G}(R/F)
}
$$
Since we also have $\pi_{F, *}\mathcal{G}(U/F) = \mathcal{G}(U/F)$
and $\pi_{F, *}\mathcal{G}(U/F) = \mathcal{G}(U/F)$
we obtain a canonical identification
$f_{small, *}\pi_{F, *}\mathcal{G}(V) = \pi_{G, *}f_*\mathcal{G}(V)$.
We omit the proof that this is compatible with restriction mappings
and that it is functorial in $\mathcal{G}$.
\end{proof}

\end{document}
